% Source PDF: source/普通高考/2012/2012陕西理.pdf, page 1, question 10.
% pdfimages -list found no embedded images; the source program flowchart is vector artwork.
% Node -> directed-edge list: 开始 -> M=0,N=0,i=1 -> random(x_i,y_i)
% -> x_i^2+y_i^2<=1; 是 -> M=M+1, 否 -> N=N+1; both -> i=i+1
% -> i>1000; 否 loops back to the random-number node, 是 -> blank process
% -> 输出P -> 结束. The blank process remains blank as in the source.
% All borders and connectors are solid; node text is centered with local parindent=0pt.

\begin{tikzpicture}[
    >=Stealth,
    x=1cm,
    y=1cm,
    line width=0.45pt,
    line cap=round,
    line join=round,
    font=\normalsize,
    flowtext/.style={anchor=center, align=center,
        execute at begin node={\setlength{\parindent}{0pt}}},
    flowbox/.style={flowtext, draw, minimum height=0.68cm,
        inner xsep=4pt, inner ysep=2pt},
    terminal/.style={flowbox, rounded corners=2pt, minimum width=1.35cm},
    process/.style={flowbox, minimum width=2.55cm},
    decision/.style={flowbox, diamond, aspect=2.8, inner sep=0pt,
        minimum width=4.25cm, minimum height=1.08cm},
    io/.style={flowbox, trapezium, trapezium left angle=72, trapezium right angle=108},
    branchlabel/.style={flowtext, inner sep=0pt},
]
    \node[terminal] (start) at (0,11.0) {开始};
    \node[process, minimum width=3.65cm] (init) at (0,9.50) {$M=0,N=0,i=1$};
    \node[process, minimum width=7.30cm] (random) at (0,8.00)
      {产生$0\sim1$之间的两个随机数分别赋给$x_i,y_i$};
    \node[decision] (inside) at (0,6.35) {$x_i^2+y_i^2\leq1$};
    \node[process, minimum width=2.25cm] (countM) at (0,4.85) {$M=M+1$};
    \node[process, minimum width=2.25cm] (countN) at (3.00,4.85) {$N=N+1$};
    \node[process, minimum width=2.10cm] (increment) at (0,3.35) {$i=i+1$};
    \node[decision, minimum width=3.75cm] (test) at (0,1.85) {$i>1000$};
    \node[process, minimum width=2.00cm] (blank) at (0,0.35) {};
    \node[io, minimum width=1.85cm] (output) at (0,-1.05) {输出 $P$};
    \node[terminal] (finish) at (0,-2.40) {结束};

    \draw[->] (start.south) -- (init.north);
    \draw[->] (init.south) -- (random.north);
    \draw[->] (random.south) -- (inside.north);

    \draw[->] (inside.south) -- node[branchlabel, right=2pt] {是} (countM.north);
    \draw (inside.east) -- (3.00,6.35) -- (3.00,5.35);
    \draw[->] (3.00,5.35) -- (countN.north);
    \node[branchlabel] at (1.35,6.70) {否};

    % The two counters merge immediately before i=i+1; only the N branch
    % carries the left-pointing arrow into the shared merge point.
    \draw[->] (countM.south) -- (0,4.10);
    \draw[->] (countN.south) -- (3.00,4.10) -- (0,4.10);
    \draw[->] (0,4.10) -- (increment.north);
    \draw[->] (increment.south) -- (test.north);

    \draw[->] (test.south) -- node[branchlabel, right=2pt] {是} (blank.north);
    \draw[->] (blank.south) -- (output.north);
    \draw[->] (output.south) -- (finish.north);

    % 否 loops left and returns with one right-pointing arrow above the
    % independent main downward arrow into random(x_i,y_i).
    \draw (test.west) -- (-3.10,1.85) -- (-3.10,8.55);
    \draw[->] (-3.10,8.55) -- (0,8.55);
    \node[branchlabel] at (-2.55,2.13) {否};

    \path[use as bounding box] (-3.30,-2.75) rectangle (4.65,11.30);
\end{tikzpicture}
